\appendix

\section{Additional Case Settings and Surrogate Validation}

\subsection{Case-setting details}

The base urban block is represented as a $240\,\mathrm{m}\times240\,\mathrm{m}$ square subdivided into a $3\times3$ grid of nine land units. Each land unit measures $80\,\mathrm{m}\times80\,\mathrm{m}$. A two-step inward offset is applied to represent the road boundary and building setback, resulting in a nominal buildable area of $60\,\mathrm{m}\times60\,\mathrm{m}$ for each unit. One land unit is designated as open space, and the remaining units are populated with residential building prototypes.

To avoid notation drift between the main text and the appendix, the following definitions are used consistently:
\begin{itemize}
\item \textbf{AF} (average floors) is defined as the total gross floor area divided by the total building footprint area.
\item \textbf{OSR} (open space ratio) is defined as the non-built-up area of the block divided by the total gross floor area.
\end{itemize}

\begin{table}[H]
\centering
\caption{Representative building prototypes used in the parametric dataset.}
\begin{threeparttable}
\begin{tabular}{@{}lcccc@{}}
\toprule
ID & Type & Footprint Area (m$^2$) & Floors & FAR Range \\
\midrule
P-1 & Point      & 2000 & 1--3   & 0.31--0.94 \\
P-2 & Point      & 1500 & 4--12  & 0.94--2.81 \\
P-3 & Point      & 1000 & 13--30 & 2.03--4.69 \\
P-4 & Point      & 1000 & 13--30 & 2.03--4.69 \\
S-1 & Slab       & 2000 & 1--3   & 0.31--0.94 \\
S-2 & Slab       & 1500 & 4--12  & 0.94--2.81 \\
S-3 & Slab       & 1000 & 13--30 & 2.03--4.69 \\
C-1 & Courtyard  & 2000 & 1--3   & 0.31--0.94 \\
C-2 & Courtyard  & 1500 & 4--12  & 0.94--2.81 \\
\bottomrule
\end{tabular}
\end{threeparttable}
\end{table}

\subsection{Surrogate validation}

The surrogate is trained on a dataset of 500 simulated morphologies. Five-fold cross-validation is used, with 400 samples for training and 100 for testing in each fold. Table~\ref{tab:supp_cv} reports the aggregate validation results corresponding to the revised surrogate assessment in the main text.

\begin{table}[H]
\centering
\caption{Five-fold cross-validation summary for the DNN surrogate model.}
\label{tab:supp_cv}
\begin{threeparttable}
\begin{tabular}{@{}lcccc@{}}
\toprule
Objective & Mean $R^2$ & Mean MAE & Std. $R^2$ & Std. MAE \\
\midrule
EUIt & 0.973 & 0.031 & $\pm 0.011$ & $\pm 0.007$ \\
EG   & 0.969 & 0.040 & $\pm 0.014$ & $\pm 0.005$ \\
H    & 0.971 & 0.028 & $\pm 0.009$ & $\pm 0.006$ \\
\bottomrule
\end{tabular}
\end{threeparttable}
\end{table}

To evaluate reliability in optimization-relevant regions, additional error slices were computed over the lower and upper tails of each target distribution. The corresponding mean absolute errors are listed in Table~\ref{tab:supp_extreme}. Table~\ref{tab:supp_fidelity} reports complementary surrogate-fidelity statistics from the cross-validation predictions.

\begin{table}[H]
\centering
\caption{Error diagnostics in extreme regions of the target distributions.}
\label{tab:supp_extreme}
\begin{threeparttable}
\begin{tabular}{@{}lccccc@{}}
\toprule
Target & MAE$_\mathrm{all}$ & MAE$_{q10}$ & MAE$_{q90}$ & $q_{0.10}$ & $q_{0.90}$ \\
\midrule
EUIt & 0.541 & 1.045 & 0.879 & 67.94 & 73.59 \\
EG   & 0.061 & 0.095 & 0.075 & 2.113 & 2.655 \\
H    & 0.075 & 0.099 & 0.077 & 6.482 & 7.513 \\
\bottomrule
\end{tabular}
\end{threeparttable}
\end{table}

\begin{table}[H]
\centering
\caption{Additional surrogate-fidelity audit on the cross-validation predictions.}
\label{tab:supp_fidelity}
\begin{threeparttable}
\begin{tabular}{@{}lccccc@{}}
\toprule
Target & RMSE & MAE & nRMSE & MAPE (\%) & $R^2$ \\
\midrule
EUIt & 0.332 & 0.254 & 0.0275 & 0.362 & 0.978 \\
EG   & 0.041 & 0.032 & 0.0385 & 1.328 & 0.962 \\
H    & 0.042 & 0.033 & 0.0234 & 0.469 & 0.987 \\
\bottomrule
\end{tabular}
\begin{tablenotes}[flushleft]
\footnotesize
\item Note: nRMSE is normalized by the observed range of each target over the cross-validation predictions.
\end{tablenotes}
\end{threeparttable}
\end{table}

\section{Benchmark Robustness and Independent Re-evaluation}

\subsection{Training dynamics and benchmark protocol}

The imported DDPG logs provide 600 episodes for each of the 20 seeds in each scenario, which makes the multi-seed learning dynamics clearer than in the earlier single-curve presentation. The revised server-side reruns used for Fig.~\ref{fig:learning} were executed under a guarded surrogate evaluator to suppress exploitation of extrapolative edge cases. Seed-level diagnostics show that the 20-episode rolling reward typically reaches 95\% of its best observed value around episodes 340 (balanced), 312 (saving), and 280 (generation). Matched-checkpoint reruns on the remote surrogate show a similar pattern, with plateau episodes around 361, 325, and 343 for the balanced, saving, and generation scenarios, respectively. However, late-stage regression remains common in both settings: the mean best-to-final reward gap ratios remain roughly 0.55--0.60 on the matched checkpoint, and 80\%--90\% of seeds still satisfy a simple substantial-regression rule.

For benchmark comparison, the revised implementation removes the previous benchmark-matching shortcut from the publication-mode NSGA-II path. Publication mode compares DDPG, random search, and NSGA-II under a shared surrogate environment and matched evaluation budget, and applies a post-hoc preference-aligned utility summary consistently across methods. On the matched remote checkpoint, DDPG remains modestly stronger than random search but weaker than fair-budget NSGA-II. The corresponding budget accounting and objective-space separation are summarized in Table~\ref{tab:supp_budget}.

\begin{table}[H]
\centering
\caption{Budget accounting and objective-space separation in the imported revision results.}
\label{tab:supp_budget}
\begin{threeparttable}
\footnotesize
\begin{tabularx}{\linewidth}{@{}>{\raggedright\arraybackslash}p{0.31\linewidth}>{\centering\arraybackslash}p{0.16\linewidth}>{\centering\arraybackslash}p{0.16\linewidth}X@{}}
\toprule
Item & Value 1 & Value 2 & Interpretation \\
\midrule
DDPG per-seed budget & 600 episodes & 40 steps/episode & 24{,}000 surrogate interactions \\
DDPG seeds per scenario & 20 & -- & pooled into one scenario archive \\
NSGA-II imported budget & 24{,}000 evaluations & 1 seed & single imported archive \\
Distance: balanced vs generation & 0.0180 & -- & normalized objective-space distance \\
Distance: balanced vs saving & 0.0145 & -- & normalized objective-space distance \\
Distance: generation vs saving & 0.0128 & -- & normalized objective-space distance \\
NSGA-II archive spread & 0.2970 & 0.1255 & mean and std of distance to normalized archive center \\
\bottomrule
\end{tabularx}
\begin{tablenotes}[flushleft]
\footnotesize
\item Note: The comparison is therefore closer to a matched per-seed budget than to a matched aggregate archive budget. The small pairwise distances among DDPG scenario means confirm that preference separation remains modest in objective space.
\end{tablenotes}
\end{threeparttable}
\end{table}

To complement the linear correlation results in the main text, Fig.~\ref{fig:supp_nonlinear} reports representative surrogate-based nonlinear response profiles for OSR$\rightarrow$EUIt, FAR$\rightarrow$EG, SVF$\rightarrow$H, and $\theta\rightarrow$H.

\begin{figure}[H]
\centering
\includegraphics[width=\linewidth]{nonlinear_response_profiles.png}
\caption{Surrogate-based nonlinear response profiles for selected UMF-target pairs.}
\label{fig:supp_nonlinear}
\end{figure}

\subsection{Independent re-evaluation of selected candidates}

To reduce over-reliance on surrogate predictions, the revised workflow re-evaluates preference-aligned candidate solutions using an independent simulation call. In the current local environment, this re-evaluation still uses the fallback analytic simulator, but the code path is structured so that a full physical-stack simulation can replace it when server-side results become available. Under the current matched-checkpoint utility selection, one DDPG candidate and two NSGA-II candidates are retained after deduplication, so Table~\ref{tab:supp_reeval} reports three unique candidates. The updated comparison shows that the selected DDPG candidate undergoes much larger EUIt drift and larger EG discrepancy than the two selected NSGA-II candidates, while its H discrepancy lies between the two NSGA-II examples.

\begin{table}[H]
\centering
\caption{Independent re-evaluation of matched-checkpoint selected candidates. Surrogate and deterministic re-evaluation values are reported for the current preference-aligned selections.}
\label{tab:supp_reeval}
\begin{threeparttable}
\scriptsize
\begin{tabular}{@{}lccccccc@{}}
\toprule
Method & Scenario & EUIt$_\mathrm{sur}$ & EUIt$_\mathrm{re}$ & EG$_\mathrm{sur}$ & EG$_\mathrm{re}$ & H$_\mathrm{sur}$ & H$_\mathrm{re}$ \\
\midrule
DDPG & Generation & 66.344 & 68.342 & 2.846 & 2.724 & 7.833 & 7.646 \\
NSGA-II & NSGA-II (seed 20260310) & 66.000 & 66.000 & 2.850 & 2.828 & 7.850 & 7.676 \\
NSGA-II & NSGA-II (seed 20260328) & 66.000 & 66.000 & 2.850 & 2.738 & 7.850 & 7.622 \\
\bottomrule
\end{tabular}
\end{threeparttable}
\end{table}
